% Companion long-record migration for Erdős #269 round 6.
% Destination label: long:torus-fourier-value
% Not frozen R. Not a blind \input. Labels may collide with the short note;
% assign destination labels as specified in R6LongRecordMigrations.md.
% Extracted from Type B edits/relocation_extracts.tex.

\subsection{An exact torus--Fourier representation of the repeated value}

The repeated three-prime value itself, rather than only its radix word, has
an explicit analytic coordinate.  For $s=2^i3^j5^k$ put
$L=\log_2s$, $\theta=\log_2 3$, and $\phi=\log_2 5$.  Let
$f_p(x)=p^{\{x\}}$ be one-periodic, with value $1$ at its jump, and let
$F_{p,M}$ be its symmetric Fourier partial sum over $|m|\le M$.

\begin{theorem}[rectangular torus--Fourier identity]
\label{res:torus-fourier-value}
The absolutely convergent sums
\[
 A_M=\sum_{i,j,k\ge0}
 \frac{F_{2,M}(L)F_{3,M}(L/\theta)F_{5,M}(L/\phi)}
      {8^i27^j125^k}
\]
satisfy
\[
 \lim_{M\to\infty}A_M
 =\sum_{i,j,k\ge0}\frac1{\hgt(2^i3^j5^k)}+\frac{17}{2}.
 \tag{9.6}\label{eq:torus-fourier-value}
\]
For each finite $M$, expansion into Fourier modes is exact.  If
$\lambda=m_1+m_2/\theta+m_3/\phi$ and
$\widehat f_p(m)=(p-1)/(\log p-2\pi i m)$, the $(m_1,m_2,m_3)$ mode is
\[
 \frac{\widehat f_2(m_1)\widehat f_3(m_2)\widehat f_5(m_3)}
 {(1-e(\lambda)/8)(1-e(\lambda\theta)/27)
  (1-e(\lambda\phi)/125)},
 \qquad e(x)=e^{2\pi i x}.
\]
\end{theorem}

\begin{proof}[Proof spine]
Floor--fractional-part decomposition gives the pointwise identity
\[
 \frac1{\hgt(s)}=
 \frac{f_2(L)f_3(L/\theta)f_5(L/\phi)}{s^3}.
\]
For fixed $M$ the mode expansion is finite, and the three lattice sums are
ordinary geometric series.  Each $f_p$ has bounded variation.  The
Dirichlet--Jordan estimate therefore bounds all $F_{p,M}$ uniformly in $M$
and gives midpoint convergence at the jump.  The summable majorant
$C\,8^{-i}27^{-j}125^{-k}$ permits dominated convergence on the exponent
lattice.

The midpoint differs from the original torus weight only on the three pure
prime axes and at the origin.  The origin contributes $8$.  At $p^n$ the
excess is $(p-1)/(2\hgt(p^n))$, while the exact jump telescope
\[
 \sum_{p\in\{2,3,5\}}(p-1)
       \sum_{n\ge1}\frac1{\hgt(p^n)}=1
\]
makes the three axes contribute $1/2$.  Their total correction is therefore
$17/2$.
\end{proof}

This is an exact ordinary Fourier-analysis theorem, not a numerical fit.
It replaces an opaque scalar value by an explicit symmetric triple-mode
limit whose frequencies are linear forms in $1,\log_3 2,\log_5 2$.
The triple Fourier series is not absolutely summable, however, and the
identity supplies no Diophantine lower bound on a linear form in the value.
The three geometric denominators are bounded away from zero by $7/8$,
$26/27$ and $124/125$; their convergence requires no small-divisor estimate
on the phase $\lambda$.  That phase is a ratio of logarithms, not
automatically an integer linear form in $\log 2,\log 3,\log 5$.  The identity
proves neither irrationality nor transcendence and carries no priority claim.
Its authored proof is
\path{ErdosProblems/Erdos269/TensorDirichletJordanCompletion.md}; the finite
Fourier computation is regression evidence rather than proof authority.
